\hypertarget{objetivos}{%
\section{Objetivos}\label{objetivos}}

\begin{frame}{Objetivos}
\protect\hypertarget{objetivos-1}{}
\begin{itemize}
\item
  Apresentar o formalismo de PEGs para descri\c{c}\~{a}o de analisadores
  sintáticos.
\item
  Apresentar o conceito de boa-forma\c{c}\~{a}o e sua rela\c{c}\~{a}o com a termina\c{c}\~{a}o
  em PEG.
\end{itemize}
\end{frame}

\hypertarget{motivauxe7uxe3o}{%
\section{Motiva\c{c}\~{a}o}\label{motivauxe7uxe3o}}

\begin{frame}{Motiva\c{c}\~{a}o}
\protect\hypertarget{motivauxe7uxe3o-1}{}
\begin{itemize}

\item
  Combinadores de parsing s\~{a}o poderosos, mas GLCs podem ser ambíguas.

  \begin{itemize}

  \item
    Como lidar com recurs\~{a}o \`{a} esquerda?
  \end{itemize}
\item
  Algumas linguagens s\~{a}o difíceis de especificar sem ambiguidade.
\end{itemize}
\end{frame}

\begin{frame}{Motiva\c{c}\~{a}o}
\protect\hypertarget{motivauxe7uxe3o-2}{}
\begin{itemize}

\item
  \textbf{PEGs}: formalismo com escolha determinística
\item
  Sem ambiguidade por defini\c{c}\~{a}o: único resultado para cada entrada
\item
  Correspondem diretamente a parsers descendentes recursivos
\end{itemize}
\end{frame}

\begin{frame}{GLC vs.~PEG: Diferen\c{c}a Fundamental}
\protect\hypertarget{glc-vs.-peg-diferenuxe7a-fundamental}{}
\begin{longtable}[]{@{}lll@{}}
\toprule\noalign{}
Aspecto & GLC & PEG \\
\midrule\noalign{}
\endhead
Escolha & N\~{a}o-determinística & \textbf{Ordenada} \\
Ambiguidade & Possível & Impossível \\
\bottomrule\noalign{}
\end{longtable}
\end{frame}

\begin{frame}{GLC vs.~PEG: Diferen\c{c}a Fundamental}
\protect\hypertarget{glc-vs.-peg-diferenuxe7a-fundamental-1}{}
\begin{longtable}[]{@{}lll@{}}
\toprule\noalign{}
Aspecto & GLC & PEG \\
\midrule\noalign{}
\endhead
Recurs\~{a}o \`{a} esquerda & Suporta naturalmente & N\~{a}o suporta (loop) \\
Predicados lookahead & N\~{a}o tem & Sim (\(!e\), \(\&e\)) \\
\bottomrule\noalign{}
\end{longtable}
\end{frame}

\begin{frame}{GLC vs.~PEG: Diferen\c{c}a Fundamental}
\protect\hypertarget{glc-vs.-peg-diferenuxe7a-fundamental-2}{}
\begin{longtable}[]{@{}lll@{}}
\toprule\noalign{}
Aspecto & GLC & PEG \\
\midrule\noalign{}
\endhead
Poder expressivo & Linguagens livres de contexto & N\~{a}o determ. \\
\bottomrule\noalign{}
\end{longtable}
\end{frame}

\begin{frame}[fragile]{Escolha N\~{a}o-Determinística vs.~Ordenada}
\protect\hypertarget{escolha-nuxe3o-determinuxedstica-vs.-ordenada}{}
\begin{itemize}

\item
  \textbf{GLC}: \(A \to \alpha \mid \beta\) --- tenta \(\alpha\)
  \textbf{e} \(\beta\) em paralelo
\item
  \textbf{PEG}: \(A \leftarrow e_1 / e_2\) --- tenta \(e_1\); \textbf{só
  se \(e_1\) falhar}, tenta \(e_2\)
\item
  A escolha ordenada \texttt{/} é determinística: no
  máximo um resultado
\item
  Essa sem\^{a}ntica elimina toda ambiguidade
\end{itemize}
\end{frame}

\hypertarget{operadores-de-peg}{%
\section{Operadores de PEG}\label{operadores-de-peg}}

\begin{frame}{Operadores de PEG}
\protect\hypertarget{operadores-de-peg-1}{}
\begin{longtable}[]{@{}lll@{}}
\toprule\noalign{}
Express\~{a}o & Nome & Sem\^{a}ntica \\
\midrule\noalign{}
\endhead
\(\varepsilon\) & vazio & Sempre tem sucesso sem consumir \\
\(a\) & terminal & Consome e casa com símbolo \(a\) \\
\(A\) & n\~{a}o-terminal & Invoca a regra de nome \(A\) \\
\bottomrule\noalign{}
\end{longtable}
\end{frame}

\begin{frame}{Operadores de PEG}
\protect\hypertarget{operadores-de-peg-2}{}
\begin{longtable}[]{@{}lll@{}}
\toprule\noalign{}
Express\~{a}o & Nome & Sem\^{a}ntica \\
\midrule\noalign{}
\endhead
\(e_1\; e_2\) & sequ\^{e}ncia & \(e_1\) depois \(e_2\) sobre o restante \\
\(e_1 / e_2\) & escolha ordenada & Tenta \(e_1\); se falhar, tenta
\(e_2\) \\
\(e^*\) & repeti\c{c}\~{a}o gulosa & Zero ou mais ocorr\^{e}ncias de \(e\) \\
\bottomrule\noalign{}
\end{longtable}
\end{frame}

\begin{frame}{Operadores de PEG}
\protect\hypertarget{operadores-de-peg-3}{}
\begin{longtable}[]{@{}lll@{}}
\toprule\noalign{}
Express\~{a}o & Nome & Sem\^{a}ntica \\
\midrule\noalign{}
\endhead
\(e^+\) & repeti\c{c}\~{a}o positiva & Uma ou mais ocorr\^{e}ncias de \(e\) \\
\(e\;?\) & opcional & Zero ou uma ocorr\^{e}ncia de \(e\) \\
\bottomrule\noalign{}
\end{longtable}
\end{frame}

\begin{frame}{Operadores de PEG}
\protect\hypertarget{operadores-de-peg-4}{}
\begin{longtable}[]{@{}
  >{\raggedright\arraybackslash}p{(\columnwidth - 4\tabcolsep) * \real{0.1364}}
  >{\raggedright\arraybackslash}p{(\columnwidth - 4\tabcolsep) * \real{0.2727}}
  >{\raggedright\arraybackslash}p{(\columnwidth - 4\tabcolsep) * \real{0.5909}}@{}}
\toprule\noalign{}
\begin{minipage}[b]{\linewidth}\raggedright
Express\~{a}o
\end{minipage} & \begin{minipage}[b]{\linewidth}\raggedright
Nome
\end{minipage} & \begin{minipage}[b]{\linewidth}\raggedright
Sem\^{a}ntica
\end{minipage} \\
\midrule\noalign{}
\endhead
\(!e\) & predicado negativo & Sucesso se \(e\) falha; n\~{a}o consome \\
\(\&e\) & predicado positivo & Sucesso se \(e\) tem sucesso; n\~{a}o
consome \\
\bottomrule\noalign{}
\end{longtable}
\end{frame}

\begin{frame}[fragile]{Escolha Ordenada}
\protect\hypertarget{escolha-ordenada}{}
\begin{lstlisting}
Se e1 tem SUCESSO:
  → Retorna o resultado de e1 imediatamente
  → e2 NUNCA é tentada

Se e1 FALHA:
  → e2 é tentada sobre a MESMA entrada original
  → Retorna o resultado de e2 (sucesso ou falha)
\end{lstlisting}

\begin{itemize}

\item
  Elimina ambiguidade: o primeiro match sempre vence
\item
  Diferen\c{c}a crucial em rela\c{c}\~{a}o \`{a} GLC
\end{itemize}
\end{frame}

\begin{frame}{Predicados: Lookahead sem Consumo}
\protect\hypertarget{predicados-lookahead-sem-consumo}{}
\begin{itemize}

\item
  \(!e\) (\textbf{predicado negativo}): tem sucesso se e somente se
  \(e\) falha; n\~{a}o consome entrada
\item
  \(\&e\) (\textbf{predicado positivo}): tem sucesso se e somente se
  \(e\) tem sucesso; n\~{a}o consome entrada
\item
  Equival\^{e}ncia: \(\&e \equiv !(! e)\)
\end{itemize}
\end{frame}

\begin{frame}[fragile]{Exemplo}
\protect\hypertarget{exemplo}{}
Exemplo --- reconhece \texttt{if} sem aceitar
\texttt{iff}:

\[\texttt{if}\; ![\text{a-zA-Z0-9}]\]

\begin{itemize}

\item
  Casa \texttt{if} (com espa\c{c}o) mas n\~{a}o
  \texttt{iff}
\item
  O predicado verifica sem consumir o caractere seguinte
\end{itemize}
\end{frame}

\begin{frame}{Repeti\c{c}\~{a}o Gulosa}
\protect\hypertarget{repetiuxe7uxe3o-gulosa}{}
A repeti\c{c}\~{a}o em PEGs é definida como:

\[e^* \leftarrow e\; e^* \;/\; \varepsilon\]
\end{frame}

\begin{frame}{Repeti\c{c}\~{a}o Gulosa}
\protect\hypertarget{repetiuxe7uxe3o-gulosa-1}{}
\begin{itemize}

\item
  A escolha é ordenada: \textbf{sempre tenta mais uma ocorr\^{e}ncia} antes
  de aceitar vazio
\item
  Resultado: \(e^*\) é sempre \textbf{guloso} --- consume o máximo
  possível
\end{itemize}
\end{frame}

\begin{frame}{Repeti\c{c}\~{a}o Gulosa}
\protect\hypertarget{repetiuxe7uxe3o-gulosa-2}{}
\begin{itemize}

\item
  Contraste com GLCs: a repeti\c{c}\~{a}o n\~{a}o é determinística
\item
  Garante que \(e^*\) sempre devolve o maior prefixo possível
\end{itemize}
\end{frame}

\hypertarget{exemplos}{%
\section{Exemplos}\label{exemplos}}

\begin{frame}{Exemplo}
\protect\hypertarget{exemplo-1}{}
\begin{itemize}

\item
  A linguagem \(\{a^nb^nc^n\,|\,n\geq 0\}\) n\~{a}o é uma LLC.

  \begin{itemize}

  \item
    Demonstra\c{c}\~{a}o usando lema do bombeamento.
  \end{itemize}
\item
  Porém, é possível construir uma PEG que a aceita.
\end{itemize}
\end{frame}

\begin{frame}{Exemplo}
\protect\hypertarget{exemplo-2}{}
\begin{itemize}

\item
  Chave: uso de predicados positivos para permitir um lookahead
  arbitrário.
\end{itemize}
\end{frame}

\begin{frame}{Exemplo}
\protect\hypertarget{exemplo-3}{}
\[
\begin{array}{lcl}
   A & \leftarrow & a A b\,/\,\epsilon\\
   B & \leftarrow & b B c\,/\,\epsilon\\
\end{array}
\]

\begin{itemize}

\item
  Express\~{a}o inicial: \(\& A a^*B\,/\,!.\)
\end{itemize}
\end{frame}

\begin{frame}{Exemplo}
\protect\hypertarget{exemplo-4}{}
\begin{itemize}

\item
  PEG para express\~{o}es
\end{itemize}

\[\begin{array}{lcl}
E &\leftarrow& T\; (\texttt{+}\; T)^* \\
T &\leftarrow& F\; (*\; F)^* \\
F &\leftarrow& \texttt{(}\; E\; \texttt{)} \;/\; [0\text{-}9]^+
\end{array}\]
\end{frame}

\begin{frame}{Exemplo}
\protect\hypertarget{exemplo-5}{}
\begin{itemize}

\item
  Associatividade \`{a} esquerda expressa por \textbf{repeti\c{c}\~{a}o}, n\~{a}o
  recurs\~{a}o
\item
  Sem recurs\~{a}o \`{a} esquerda (problema de parsers descendentes)
\item
  A GLC equivalente precisaria de recurs\~{a}o \`{a} esquerda ou transforma\c{c}\~{a}o
\end{itemize}
\end{frame}

\begin{frame}{Compara\c{c}\~{a}o}
\protect\hypertarget{comparauxe7uxe3o}{}
GLC com preced\^{e}ncia: \[E \to E + T \mid T \qquad T \to T * F \mid F\]
\end{frame}

\begin{frame}{Compara\c{c}\~{a}o}
\protect\hypertarget{comparauxe7uxe3o-1}{}
PEG equivalente:
\[E \leftarrow T\; (\texttt{+}\; T)^* \qquad T \leftarrow F\; (*\; F)^*\]

\begin{itemize}

\item
  PEG é mais direta: n\~{a}o precisa de recurs\~{a}o \`{a} esquerda
\item
  O operador \((\cdot)^*\) captura a repeti\c{c}\~{a}o iterativamente
\end{itemize}
\end{frame}

\hypertarget{terminauxe7uxe3o-e-boa-formauxe7uxe3o}{%
\section{Termina\c{c}\~{a}o e
Boa-Forma\c{c}\~{a}o}\label{terminauxe7uxe3o-e-boa-formauxe7uxe3o}}

\begin{frame}{Termina\c{c}\~{a}o}
\protect\hypertarget{terminauxe7uxe3o}{}
Uma PEG é \textbf{bem formada} (e portanto termina) se:

\begin{enumerate}

\item
  \textbf{N\~{a}o há recurs\~{a}o \`{a} esquerda} direta ou indireta
\item
  \textbf{Toda repeti\c{c}\~{a}o \(e^*\) é tal que \(e\) consume ao menos um
  símbolo} (n\~{a}o pode ter sucesso sem consumir entrada)
\end{enumerate}
\end{frame}

\begin{frame}{Termina\c{c}\~{a}o}
\protect\hypertarget{terminauxe7uxe3o-1}{}
\begin{itemize}

\item
  Condi\c{c}\~{a}o 1: recurs\~{a}o \`{a} esquerda causaria loop infinito em parsers
  descendentes
\item
  Condi\c{c}\~{a}o 2: se \(e\) pode ter sucesso sem consumir, \(e^*\) repetiria
  indefinidamente
\end{itemize}
\end{frame}

\hypertarget{implementauxe7uxe3o-em-haskell}{%
\section{Implementa\c{c}\~{a}o em
Haskell}\label{implementauxe7uxe3o-em-haskell}}

\begin{frame}[fragile]{O Tipo \texttt{Result}}
\protect\hypertarget{o-tipo-result}{}
\begin{lstlisting}[language=Haskell]
data Result d a
  = Pure a           -- sucesso sem consumir entrada
  | Commit d a       -- sucesso consumindo entrada (d = resto)
  | Fail String Bool -- falha (Bool: houve consumo?)
\end{lstlisting}
\end{frame}

\begin{frame}[fragile]{O Tipo \texttt{Result}}
\protect\hypertarget{o-tipo-result-1}{}
\begin{itemize}

\item
  \texttt{Pure}: sucesso sem consumo → pode tentar
  alternativa
\item
  \texttt{Commit}: sucesso com consumo → sem retrocesso
  possível
\end{itemize}
\end{frame}

\begin{frame}[fragile]{O Tipo \texttt{Result}}
\protect\hypertarget{o-tipo-result-2}{}
\begin{itemize}

\item
  \texttt{Fail \_ False}: falha sem consumo → pode
  tentar alternativa
\item
  \texttt{Fail \_ True}: falha \textbf{com consumo} →
  \textbf{sem} retrocesso
\end{itemize}
\end{frame}

\begin{frame}[fragile]{O Tipo \texttt{PExp}}
\protect\hypertarget{o-tipo-pexp}{}
\begin{lstlisting}[language=Haskell]
newtype PExp d a = PExp { runPExp :: d -> Result d a }
\end{lstlisting}

\begin{itemize}

\item
  \texttt{d}: tipo da entrada (tipicamente
  \texttt{String})
\item
  \texttt{a}: tipo do resultado
\item
  Inst\^{a}ncia de \texttt{Functor},
  \texttt{Applicative},
  \texttt{Alternative}, \texttt{Monad}
\end{itemize}
\end{frame}

\begin{frame}[fragile]{Sequ\^{e}ncia}
\protect\hypertarget{sequuxeancia}{}
\begin{lstlisting}[language=Haskell]
instance Applicative (PExp d) where
  pure a = PExp $ \_ -> Pure a
  PExp mf <*> PExp ma = PExp $ \d ->
    case mf d of
      Pure f      -> fmap f (ma d)
      Fail s c    -> Fail s c
      Commit d' f ->
        case ma d' of
          Pure a       -> Commit d' (f a)
          Fail s _     -> Fail s True  -- já consumiu!
          Commit d'' a -> Commit d'' (f a)
\end{lstlisting}
\end{frame}

\begin{frame}[fragile]{Escolha}
\protect\hypertarget{escolha}{}
\begin{lstlisting}[language=Haskell]
instance Alternative (PExp d) where
  PExp ma <|> PExp mb = PExp $ \d ->
    case ma d of
      Fail _ False -> mb d  -- falhou SEM consumir: tenta alternativa
      x            -> x     -- sucesso OU falhou COM consumo: retorna
  empty = PExp $ \_ -> Fail "empty" False
\end{lstlisting}
\end{frame}

\begin{frame}[fragile]{O Combinador \texttt{try}}
\protect\hypertarget{o-combinador-try}{}
\begin{lstlisting}[language=Haskell]
try :: PExp d a -> PExp d a
try (PExp m) = PExp $ \d ->
  case m d of
    Fail s _ -> Fail s False  -- converte qualquer falha em "sem consumo"
    x        -> x
\end{lstlisting}
\end{frame}

\begin{frame}[fragile]{O Combinador \texttt{try}}
\protect\hypertarget{o-combinador-try-1}{}
\begin{itemize}

\item
  \texttt{try p} executa \texttt{p}
  normalmente
\item
  Em caso de falha, \textbf{cancela} o consumo para fins de backtracking
\item
  Permite retroceder mesmo que \texttt{p} tenha
  consumido parte da entrada
\end{itemize}
\end{frame}

\begin{frame}[fragile]{Escolha Ordenada}
\protect\hypertarget{escolha-ordenada-1}{}
\begin{lstlisting}[language=Haskell]
infixl 3 </>
(</>) :: PExp d a -> PExp d a -> PExp d a
p </> q = try p <|> q
\end{lstlisting}
\end{frame}

\begin{frame}[fragile]{Escolha Ordenada}
\protect\hypertarget{escolha-ordenada-2}{}
\begin{itemize}

\item
  \texttt{p </> q}: tenta \texttt{p}
  com backtracking garantido
\item
  Se \texttt{p} falha por qualquer raz\~{a}o,
  \texttt{q} é tentado sobre a entrada original
\item
  Implementa diretamente a sem\^{a}ntica de PEGs
\end{itemize}
\end{frame}

\begin{frame}[fragile]{O Predicado Negativo
\texttt{not}}
\protect\hypertarget{o-predicado-negativo-not}{}
\begin{lstlisting}[language=Haskell]
not :: PExp d a -> PExp d ()
not (PExp m) = PExp $ \d ->
  case m d of
    Fail{} -> Pure ()            -- e falhou: sucesso sem consumir
    _      -> Fail "unexpected" False  -- e teve sucesso: falha
\end{lstlisting}
\end{frame}

\begin{frame}[fragile]{Exemplo}
\protect\hypertarget{exemplo-6}{}
\begin{lstlisting}[language=Haskell]
eof :: Stream d => PExp d ()
eof = not anyChar  -- sucesso apenas no fim da entrada
\end{lstlisting}
\end{frame}

\begin{frame}[fragile]{Combinadores Léxicos}
\protect\hypertarget{combinadores-luxe9xicos}{}
\begin{lstlisting}[language=Haskell]
satisfy :: Stream d => (Char -> Bool) -> PExp d Char
satisfy p = try $ do
  x <- anyChar
  x <$ guard (p x)
\end{lstlisting}
\end{frame}

\begin{frame}[fragile]{Combinadores Léxicos}
\protect\hypertarget{combinadores-luxe9xicos-1}{}
\begin{lstlisting}[language=Haskell]
lexeme :: Stream d => PExp d a -> PExp d a
lexeme m = m <* whiteSpace
\end{lstlisting}
\end{frame}

\begin{frame}[fragile]{Combinadores léxicos}
\protect\hypertarget{combinadores-luxe9xicos-2}{}
\begin{lstlisting}[language=Haskell]
symbol :: Stream d => Char -> PExp d Char
symbol c = lexeme (char c)
\end{lstlisting}
\end{frame}

\begin{frame}[fragile]{Parser de Express\~{o}es}
\protect\hypertarget{parser-de-expressuxf5es}{}
\begin{lstlisting}[language=Haskell]
data Exp = Lit Int | Exp :+: Exp | Exp :*: Exp

expP :: Stream d => PExp d Exp
expP = foldl (:+:) <$> termP <*> many (symbol '+' *> termP)
\end{lstlisting}
\end{frame}

\begin{frame}[fragile]{Parser de Express\~{o}es}
\protect\hypertarget{parser-de-expressuxf5es-1}{}
\begin{lstlisting}[language=Haskell]
termP :: Stream d => PExp d Exp
termP = foldl (:*:) <$> factP <*> many (symbol '*' *> factP)
\end{lstlisting}
\end{frame}

\begin{frame}[fragile]{Parser de Express\~{o}es}
\protect\hypertarget{parser-de-expressuxf5es-2}{}
\begin{lstlisting}[language=Haskell]
factP :: Stream d => PExp d Exp
factP = between (symbol '(') (symbol ')') expP
    </> Lit . foldl (\a b -> a*10+b) 0 <$> many1 digit
\end{lstlisting}
\end{frame}

\hypertarget{conclusuxe3o}{%
\section{Conclus\~{a}o}\label{conclusuxe3o}}

\begin{frame}{Conclus\~{a}o}
\protect\hypertarget{conclusuxe3o-1}{}
\begin{itemize}

\item
  PEGs usam \textbf{escolha ordenada} \(e_1 / e_2\) em vez de
  n\~{a}o-determinística
\item
  S\~{a}o \textbf{n\~{a}o-ambíguas por defini\c{c}\~{a}o}: o primeiro match sempre vence
\end{itemize}
\end{frame}

\begin{frame}{Conclus\~{a}o}
\protect\hypertarget{conclusuxe3o-2}{}
\begin{itemize}

\item
  \textbf{Predicados} \(!e\) e \(\&e\) permitem lookahead de tamanho
  arbitrário
\item
  A \textbf{repeti\c{c}\~{a}o} \(e^*\) é gulosa por natureza da escolha ordenada
\end{itemize}
\end{frame}

\begin{frame}[fragile]{Conclus\~{a}o}
\protect\hypertarget{conclusuxe3o-3}{}
\begin{itemize}

\item
  PEGs bem formadas terminam sempre (sem recurs\~{a}o \`{a} esquerda, sem ciclos
  vazios)
\item
  Implementa\c{c}\~{a}o: o tipo \texttt{Result} distingue falha
  com e sem consumo
\item
  \texttt{try} e \texttt{</>}
  implementam backtracking explícito
\end{itemize}
\end{frame}
